\documentclass[10pt,a4paper]{amsart}
\usepackage[margin=19mm]{geometry}
\usepackage{amsmath,amssymb,mathtools}
\usepackage[scr=rsfs,cal=euler]{mathalfa}
\usepackage{xfrac}
\usepackage[unicode,hidelinks,bookmarks=false]{hyperref}
\newcommand{\ie}{i.e.\ }
\newcommand{\cf}{cf.\ }
\newcommand{\resp}{respectively\ }
\newcommand{\Subsection}{\subsection}
\providecommand{\nobreakdash}{\nobreak\hbox{-}}
\setlength{\emergencystretch}{2em}
\multlinegap0pt
\usepackage{bm}
\usepackage{bbm}
\usepackage{dsfont}
\usepackage{xspace}
\usepackage{cancel}
\usepackage{mathtools}
\usepackage{amsthm}
\makeatletter
\@ifundefined{proposition}{\newtheorem{proposition}{Proposition}}{}
\@ifundefined{remark}{\newtheorem{remark}{Remark}}{}
\makeatother
\newcommand{\dd}{{\, \mathrm d}}
\title[Global Well-Posedness near Rayleigh-Jeans Equilibria for the]{Global Well-Posedness near Rayleigh-Jeans Equilibria for the Cubic NLS Wave Kinetic Equation}
\author{Miguel Escobedo, Angeliki Menegaki}
\date{Source: arXiv:2608.19483v1 (CC BY 4.0)}
\begin{document}
\maketitle

\noindent\emph{Source and licence.} Global Well-Posedness near Rayleigh-Jeans Equilibria for the Cubic NLS Wave Kinetic Equation. Version arXiv:2608.19483v1 (CC BY 4.0), distributed under the Creative Commons Attribution 4.0 International licence, as recorded in the arXiv licence metadata for this exact version and shown on the abstract page https://arxiv.org/abs/2608.19483v1. Full rights notice: licenses/arxiv-2608.19483-CC-BY-4.0.txt.\par\smallskip

\noindent\textit{Editorial scope.} Counted unit: the Remark whose source label is \texttt{None} in the article (source0.tex lines 1099--1104, source label \texttt{None}), delivered together with the article's own complete original proof of it (source0.tex lines 1105--1573, one \texttt{proof} environment of 469 lines). No mathematics is written, repaired or re-derived: statement and proof are the article's own text line for line. Required context retained uncounted: prop:noncompactness (lines 711--713). Every definition, display and label of those lines is retained unchanged. Omitted from this package and not redistributed: 1--710, 714--1098, 1574--2406. Nothing inside a retained excerpt is omitted or altered: every retained body is delivered exactly as the article's lines are written, so no omission receipt of any kind accompanies this package. Citations: the retained excerpts carry no citation command at all, so no bibliography entry of the article is transcribed and no reference is redistributed. Reference adaptation: none was needed and none was made; every reference inside the delivered unit resolves inside it, so no unresolved reference or citation is produced. Printed slips kept literal: no printed word, symbol, hypothesis or number of the counted statement or of its proof was altered. Layout and class: the article's own class is \texttt{amsart} with packages including amsmath, amssymb, amsthm, graphicx, mathrsfs, url, color, hyperref, amsxtra, bbm, epstopdf, xcolor, verbatim, geometry; the wrapper is the lane's pinned \texttt{amsart} editorial wrapper at 10pt on A4, which restates the theorem-like environments the retained text needs and adds the article's own macro definitions that the retained text needs (\texttt{\textbackslash dd}), with extra wrapper packages (bm, bbm, dsfont, xspace, cancel, mathtools). The delivered numbering is the unit's own. Assets: no retained span contains a figure environment, a graphics-inclusion command, a PDF-page inclusion, a table or a TikZ picture; no image file is copied into this package, the package declares no asset dependency and no image file is redistributed with it. Licence: version arXiv:2608.19483v1 of this article is distributed under the Creative Commons Attribution 4.0 International licence, as recorded in the arXiv licence metadata for this exact version and shown on the abstract page https://arxiv.org/abs/2608.19483v1. A full rights notice is delivered at licenses/arxiv-2608.19483-CC-BY-4.0.txt. Characters: every retained excerpt is pure ASCII, so the delivered engine, XeLaTeX with its default Latin Modern text font, reports no missing character.
\begin{proposition} \label{prop:noncompactness}
  For all $\mu  \geq 0$, the operator $K_\mu$ is not compact on $L^2(\sqrt{\omega} \dd \omega)$.
\end{proposition}

\begin{remark}[On the dependency of the constant $C_\mu$ on $\mu$]
Note that the estimate of $C_\mu $ above depends on $\mu$ and in particular it blows up as $\mu \to 0$. Indeed using the same dilation unitary operator as in the proof of Prop. \ref{prop:noncompactness}, $(U_\lambda g)(\omega) := \lambda^{-3/4} g(\lambda^{-1}\omega)$, we see that for example for the quadratic nonlinearities: $U_\lambda^{-1} \Gamma_{\mu} (U_\lambda g,U_\lambda g) = \lambda^{-3/4}\Gamma_{\tfrac{\mu}{\lambda}} (g,g) $. Indeed briefly, $(U_\lambda^{-1}g)(\omega)=\lambda^{3/4}g(\lambda\omega)$ and setting $ \omega_i=\lambda x_i$ in the integral, the measure contributes $\lambda^2$, the two RJ  factors contribute $\lambda^{-2}$, since $$ n_\mu(\lambda x)=\lambda^{-1}n_{\mu/\lambda}(x),  $$ the two factors $U_\lambda g$ contribute $\lambda^{-3/2}$, while the kernel is 0-homogeneous, and $U_\lambda^{-1}$ contributes $\lambda^{3/4}$. Altogether, $$ U_\lambda^{-1}\Gamma_\mu(U_\lambda g,U_\lambda g) = \lambda^{3/4+2-2-3/2} \Gamma_{\mu/\lambda}(g,g) = \lambda^{-3/4}\Gamma_{\mu/\lambda}(g,g). $$
Then $$\| U_\lambda^{-1} \Gamma_{\mu} (U_\lambda g,U_\lambda g)\|_{L^2(\sqrt{\omega})} = \| \Gamma_{\mu} (U_\lambda g,U_\lambda g) \|_{L^2(\sqrt{\omega})}  = \lambda^{-3/4} \|\Gamma_{\tfrac{\mu}{\lambda}} ( g,g) \|_{L^2(\sqrt{\omega})}. $$
For $\mu=\lambda$, we obtain the bound $C_1\mu^{-3/4}$. For the cubic term, a similar scaling in $\mu$ holds with power $-3/2$ instead of $-3/4$. Thus indeed our hypothesis $\mu>0$ is not just a technicality. 
%shows that our bounds of the nonlinear terms blow up as $\mu \to 0$. 
\end{remark}

    \begin{proof}
        \noindent
\underline{\emph{We start with the cubic nonlinearities}} which are more difficult. 
        We first notice that the terms for which the nonlinearity is of the form $$n_i^{-1}g_0g_kg_\ell \quad \text{ for }\ i, k, \ell \in \{1,2,3\},$$
        i.e. have $g_0$ as a factor, can be treated as follows. For example for the first term we write: 
        \begin{align*}
            &\left\| \iint_{\substack{\omega_1, \omega_2>0 \\ \omega_1+\omega_2>\omega_0}}  \frac{\text{min}(\sqrt{\omega_0},\sqrt{\omega_1},\sqrt{\omega_2},\sqrt{\omega_1+\omega_2-\omega_0})}{\sqrt{\omega_0}} 
   n_1n_2 g_0g_1g_2 \right\|_{L^2(\sqrt{\omega_0})}^2 \\
   & \leq \int_0^\infty \sqrt{\omega_0} |g_0|^2 \left\vert 
   \iint_{\substack{\omega_1, \omega_2>0 \\ \omega_1+\omega_2>\omega_0}}  \frac{\text{min}(\sqrt{\omega_0},\sqrt{\omega_1},\sqrt{\omega_2},\sqrt{\omega_3})}{\sqrt{\omega_0}} n_1n_2 g_1g_2 \dd \omega_1 \dd \omega_2
   \right\vert^2 d \omega  \\
   & \leq \|g\|_{L^2(\sqrt{\omega} d\omega)}^2\|g\|_{L^2(\sqrt{\omega} d\omega)}^4 = \|g\|_{L^2(\sqrt{\omega} d\omega)}^6, \end{align*}
        after bounding the minimum factor by $1$, multiplying and dividing by $\omega^{1/4}$ and applying H\"{o}lder's inequality in the inner double integral. The remaining two terms that look like this, are quite similar and thus omitted. 
        
Therefore,  the only delicate term in the cubic nonlinearity is the one with the nonlinearity of the form $n_0^{-1} g_1g_2g_3$. We call this integral $T_{\text{del}}(g,g,g)$.  
To bound this term, the idea is to proceed by a dyadic decomposition of
our integrals. We split the integral into annuli for $j \geq 0$,
so that the integral can be decomposed into 
\begin{align*} 
 I_j &:= \{ \omega >0: 2^{j} < \omega \leq 2^{j+1}\}, \quad \text{for } j\geq 0,  \\
 I_{-1}&:= \{ \omega >0: \omega \in (0,1]\}.
\end{align*}
We then control each integral on these annuli. 

 We start by writing $g = g 1_{I_{-1}} + \sum_{j \geq 0} g 1_{I_j}$ and assuming that $\omega_i \in I_{j_i}$ for each $i=-1,1,2,3,0$. We first going to show the estimate for $\min(j_0, j_1,j_2,j_3) \geq 0$, that is when all $\omega_{i}$'s are larger than $1$. 
 The following concerns bounding the localised-for-each-frequency delicate term  
        $$ 1_{I_{j_0}} T_{\text{del}}(g_{j_1},g_{j_2},g_{j_3}), \quad 
        \text{ where }\  g_{j_i} = g  1_{I_{j_i}}, 
        $$
        since the original term is recovered by gluing the localised terms together:
        \begin{align*}
             \|T_{\text{del}}(g,g,g)\|_{L^2(\sqrt{\omega})}& = 
            \left[\int_0^\infty \sqrt{\omega} |T_{\text{del}}|^2 \dd \omega \right]^{1/2}\\
            &= 
            \left[ \sum_{j_0\geq 0} \int_{I_{j_0}} \sqrt{\omega_0} \left\vert  
            \sum_{j_1, j_2, j_3} T_{\text{del}}(g_{j_1},g_{j_2},g_{j_3})
            \right\vert^2 \dd \omega_0 \right]^{1/2} \\
            &  = \left[\sum_{j_0\geq 0}
            \left\| \sum_{j_1, j_2, j_3} T_{\text{del}}(g_{j_1},g_{j_2},g_{j_3})
            \right\|_{L^2(I_{j_0};\sqrt{\omega}\dd \omega)}^2   \right]^{1/2} \\
            & \leq 
            \left[\sum_{j_0\geq 0} \left( 
            \sum_{j_1,j_2,j_3}
            \|T_{\text{del}}(g_{j_1},g_{j_2},g_{j_3})  \|_{L^2(I_{j_0};\sqrt{\omega}\dd \omega)} \right)^2 
            \right]^{1/2}
        \end{align*}
        
        where the last step is just the triangle inequality. 
        Above we may also assume that $j_1\geq j_2 \geq j_3$. We can do that up to a factor $6$ of all the possible permutations, since $\sum_{j_1,j_2,j_3} = \sum_{\text{permutations}} \sum_{j_{\sigma(1)} \geq j_{\sigma(2)} \geq j_{\sigma(3)}}$.  

        We estimate now the integral kernel:
        \begin{align}
            \sigma_{0,1,2,3}:&= \frac{\text{min}(\sqrt{\omega_0},\sqrt{\omega_1},\sqrt{\omega_2},\sqrt{\omega_3})}{\sqrt{\omega_0}} 
   \frac{n_1n_2n_3}{n_0}\notag \\
   &\lesssim \min \left(1,
   \sqrt{
   \frac{2^{\min(j_1,j_2,j_3)}}{2^{j_0}} 
   } \right) 2^{j_0 - j_1-j_2-j_3}  = \min \left(1,
   \sqrt{
   \frac{2^{j_3}}{2^{j_0}} 
   } \right) 2^{j_0 - j_1-j_2-j_3}
        \end{align}
        because we have assumed that $j_3$ is the minimum scale among the $j_1,j_2,j_3$. So then 
        \begin{equation} \label{eq:cubic term_dydadic}
            \begin{split}
                &\|T_{\text{del}}(g_{j_1},g_{j_2},g_{j_3})  \|_{L^2(I_{j_0};\sqrt{\omega}\dd \omega)}   = 
                \left[ \int_{I_{j_0}}\sqrt{\omega_0} \left\vert
                \iint \sigma_{0,1,2,3}  \ g_{j_1}g_{j_2}g_{j_3} \dd \omega_1 \dd \omega_2 
                \right\vert^2 \dd \omega_0\right]^{1/2}\\
                & \lesssim \min \left(1, \sqrt{
     \frac{2^{j_3}}{2^{j_0}} 
   } \right) 2^{j_0-j_1-j_2-j_3} 2^{j_0/4} 
   \left[ \int_{I_{j_0}}  \left\vert 
   \iint g_{j_1}g_{j_2}g_{j_3}\dd \omega_1\dd \omega_2
   \right\vert^2
   \right]^{1/2} \\
   &  = \min \left(1, 2^{(j_3-j_0)/2} \right) 2^{5 j_0/4-j_1-j_2-j_3}
   \| g_{j_1} \ast g_{j_2} \ast \overline{g_{j_3}} \|_{L^2(\dd \omega)}
            \end{split}
        \end{equation}
        where the factor $2^{j_0/4}$ comes due to the weight $\sqrt{\omega_0}$. Here we extend each $g_{j_3}$ by zero to the whole real line and denote by
$\overline{g_{j_3}}(z):=g_{j_3}(-z)$ its reflection. With this convention,
the resonance relation $\omega_3=\omega_1+\omega_2-\omega_0$ allows us to write the integral as a convolution on $\mathbb{R}$. This is necessary as the convolution makes $g_{j_3}$ to be valued at $-\omega_3$. 
        
        We continue by applying Young's Inequality to the convolution term in \eqref{eq:cubic term_dydadic}: 
        \begin{equation}\label{eq:Young_Ineq_convolution}
         \| g_{j_1} \ast g_{j_2} \ast \overline{g_{j_3}} \|_{L^2(\dd \omega)} \leq \|g_{j_1} \|_{L^2(\dd \omega)} \|g_{j_2} \|_{L^1(\dd \omega)} \| g_{j_3} \|_{L^1(\dd \omega)}, 
         \end{equation}
        where we choose to place the $L^2$ norm in the term localised in the highest frequency (here assumed to be $j_1$) so that we extract more decay. To get the estimate in our weighted topology $L^2(\sqrt{\omega}\dd \omega)$, we compute explicitly 
        \begin{align} \label{eq:estimL2_1}
            & \|g_{j_i}\|_{L^2(\dd \omega)} \lesssim 2^{-j_i/4} \|g_{j_i}\|_{L^2(\sqrt{\omega}\dd\omega)}\\
            \label{eq:estimL2_2}
            & \|g_{j_i}\|_{L^1(\dd \omega)} \lesssim 
            2^{j_i/4} \|g_{j_i}\|_{L^2(\sqrt{\omega}\dd\omega)}.
        \end{align}
        Inserting these estimates in \eqref{eq:cubic term_dydadic} we get  
        \begin{equation}
            \begin{split}
                 \|T_{\text{del}}(g_{j_1},g_{j_2},g_{j_3})  \|_{L^2(I_{j_0};\sqrt{\omega}\dd \omega)}  
                 \lesssim  
                &\min \left(1, 2^{(j_3-j_0)/2} \right) 2^{5 j_0/4 - 5j_1/4 - 3j_2/4 - 3j_3/4} \times  \\ &\times \|g_{j_1}\|_{L^2(\sqrt{\omega}\dd\omega)}
                \| g_{j_2}\|_{L^2(\sqrt{\omega}\dd\omega)}
                \| g_{j_3}\|_{L^2(\sqrt{\omega}\dd\omega)}.
            \end{split}
        \end{equation}
        Therefore, altogether we get 
        \begin{equation}
            \begin{split}
                        &\|T_{\text{del}}(g,g,g)\|_{L^2(\sqrt{\omega})} 
                        \lesssim  
                       \Bigg[ \sum_{j_0\geq 0}
                       \Bigg( \sum_{j_1\geq j_2 \geq j_3}   
                       \min \left(1, 2^{(j_3-j_0)/2} \right) \times \\ 
                       &  \times 
                    2^{5 j_0/4 - 5j_1/4 - 3j_2/4 - 3j_3/4}\|g_{j_1}\|_{L^2(\sqrt{\omega}\dd\omega)}
                \|g_{j_2}\|_{L^2(\sqrt{\omega}\dd\omega)}
                \| g_{j_3}\|_{L^2(\sqrt{\omega}\dd\omega)} \Bigg)^2 \Bigg]^{1/2}.
            \end{split}
        \end{equation}
        
        Now we split into two possible cases on whether $j_0$ is larger or smaller to $j_1$. \\

        \noindent
        \underline{Case 1: If $j_0 \geq j_1 - \text{Const.}$:} The resonant relation forces then that 
        $j_0 \sim j_1$. Indeed we have 
         $2^{j_0}\sim \omega_0  = \omega_1+\omega_2 - \omega_3 \leq \omega_1+\omega_2 \leq 2^{j_1+1}$ since $j_1$ is assumed to be the largest scale. So if 
$j_0 \geq j_1 - \text{Const.}$, it has to be the case that $j_0 \sim j_1$. In other words, $\omega_1$ and $\omega_0$ live on comparable dyadic scales, in that scenario.
%j_1-constant \leq j_0 \leq j_1+1
Inserting that in our main term, we don't have the sum over $j_1$ anymore (since $j_1$ and $j_0$ differ by at most a fixed constant). Then 
\begin{equation}
    \begin{split}
      &\|T_{\text{del}}(g,g,g)\|_{L^2(\sqrt{\omega})} 
        \lesssim \\ & 
        \left[\sum_{j_0\geq 0}
                       \left( \sum_{j_0\geq j_2 \geq j_3} 
                       2^{(j_3-j_0)/2}
                       2^{-3j_2/4 - 3j_3/4}
\|g_{j_0}\|_{L^2(\sqrt{\omega}\dd\omega)}
 \| g_{j_2}\|_{L^2(\sqrt{\omega}\dd\omega)}
\| g_{j_3}\|_{L^2(\sqrt{\omega}\dd\omega)}\right)^2 \right]^{1/2}\\
& = \left[\sum_{j_0\geq 0}
2^{-j_0}
\|g_{j_0}\|_{L^2(\sqrt{\omega}\dd\omega)}^2 
              \left( \sum_{j_0\geq j_2}       
 2^{-3j_2/4} \| g_{j_2} \|_{L^2(\sqrt{\omega}\dd\omega)}
 \sum_{j_3 \leq j_2 } 2^{-j_3/4} 
\| g_{j_3}\|_{L^2(\sqrt{\omega}\dd\omega)}\right)^2 \right]^{1/2}
\\
    & \lesssim \left[\sum_{j_0\geq 0} 2^{-j_0}\|g_{j_0}\|_{L^2(\sqrt{\omega}\dd\omega)}^2
                \right]^{1/2} \|g\|_{L^2(\sqrt{\omega}\dd\omega)}^2 \lesssim \|g\|_{L^2(\sqrt{\omega}\dd\omega)}^3.
           \end{split}
        \end{equation}
where in the last inequality we applied Cauchy-Schwarz. 

        \noindent
        \underline{Case 2: If $j_0 < j_1 - \text{Const.}$:} In this case the resonant condition forces $j_1\sim j_2 \sim j_3$. Indeed the resonances give that $\omega_1+ \omega_2=\omega_0+\omega_3$ and the left-hand side is dominated by $\omega_1$ since $j_1$ is the largest scale and $\omega_1+\omega_2 \sim 2^{j_1}$. In order to balance the large $\omega_1$ on the left, the right-side must contain a comparable large term. But  given that $\omega_0< \omega_1$, it must be that $\omega_3 \sim \omega_1$ or equivalently that $j_3 \sim j_1$. Now, since we have assumed that $j_1\geq j_2\geq j_3$ it has to be that all  frequencies are in comparable dyadic scales and thus $j_1\sim j_2 \sim j_3$. 

        Inserting that in our main term, we just have two main sums, since $j_1,j_2,j_3$ differ by at most constants. So:
        \begin{equation}
    \begin{split}
      &\|T_{\text{del}}(g,g,g)\|_{L^2(\sqrt{\omega})} 
        \lesssim \\ & 
        \left[ \sum_{j_0\geq 0}
                       \left( \sum_{j_1> j_0 + \text{Const.}}   
                        2^{5 j_0/4 - 11 j_1/4}
            \|g_{j_1}\|_{L^2(\sqrt{\omega}\dd\omega)}^3
                 \right)^2 \right]^{1/2} \\
                 & \lesssim 
                 \sum_{j_1\geq 0} 
                 \left[ \sum_{j_0} \left(1_{j_1>j_0+\text{Const.}} 2^{5 j_0/4 - 11 j_1/4}\|g_{j_1}\|_{L^2(\sqrt{\omega}\dd\omega)}^3\right)^2
                 \right]^{1/2}\\
                 & = 
                 \sum_{j_1\geq 0} 2^{- 11 j_1/4} \|g_{j_1}\|_{L^2(\sqrt{\omega}\dd\omega)}^3 
                 \left( \sum_{j_0< j_1 - \text{Const.}} 2^{5 j_0/2}  \right)^{1/2}\\
                 & \lesssim \sum_{j_1\geq 0} 2^{- 11 j_1/4} \|g_{j_1}\|_{L^2(\sqrt{\omega}\dd\omega)}^3 2^{5j_1/4}
                 \\
                 & = 
                 \sum_{j_1\geq 0} 2^{- 3 j_1/2} \|g_{j_1}\|_{L^2(\sqrt{\omega}\dd\omega)}^3  \lesssim \|g\|_{L^2(\sqrt{\omega}\dd\omega)}^3 
        \end{split}
        \end{equation}
        where the minimum-factor in the beginning is just $1$, for the second inequality we applied triangle inequality in the $\ell^2_{j_0}$-norm, and from the fourth to fifth line we bounded the geometric series by $2^{5(j_1 -\text{Const.})/2}\lesssim 2^{5 j_1/2}$.
        
For the remaining cases regarding the ordering of the scales $j_1, j_2, j_3$, i.e. when $j_3$ or $j_2 = \max(j_1,j_2, j_3)$, the argument is very similar and thus omitted. What is important in each case is to share the $L^2-L^1-L^1$ norm in the Young's inequality for convolutions, in \eqref{eq:Young_Ineq_convolution},  by placing the $L^2$ norm to the highest scales assumed. 

Finally notice that so far we treated all the scenaria as long as all the frequencies are larger than $1$. This suffices since the low-frequency block $I_{-1}$ is harmless for $\mu>0$: Indeed, we can use that on bounded frequencies, $n_i=(\omega+\mu)^{-1}$ and $n_i^{-1}$ are all bounded and estimates like \eqref{eq:estimL2_2} are still true: 
$$
\|g_{-1}\|_{L^1(\dd\omega)} \leq \left(\int_0^1 |g(\omega)|^2\sqrt{\omega} \dd \omega\right)^{1/2}
\left(\int_0^1 \omega^{-1/2} \dd\omega\right)^{1/2}
\lesssim
\|g_{-1}\|_{L^2(\sqrt{\omega}\dd\omega)}.
$$
In the Young convolution estimates above, the $L^2(\dd\omega)$ factor is always chosen to be in the highest frequency term.\\
So, \emph{if the highest frequency is not in the low block $I_{-1}$}, we use
\eqref{eq:estimL2_1} for that highest-frequency factor, 
while the remaining lower frequency factors are placed in $L^1(\dd\omega)$ slot as was done above. \\

\emph{If the highest frequency belongs in $I_{-1}$}, then all input frequencies are bounded: indeed
 if we have $j_1\geq j_2\geq j_3$ and $j_1=-1$ then $j_2= j_3=-1$ as well. Also by the resonance relation $\omega_0= \omega_1+\omega_2-\omega_3<2$, and thus bounded, too.
 In this case the space  $L^2((0,1);\sqrt{\omega}\dd\omega)$ is not a subspace of $L^2((0,1);\dd\omega)$ so we need to adjust the estimate \eqref{eq:estimL2_1} when applying Young's inequality for convolutions.
In fact we have $L^2((0,1);\sqrt{\omega}\dd\omega) \subset L^{\frac{6}{5}}((0,1))$: Indeed by H\"{o}lder's  
\begin{align*}
\int_0^1 |g(\omega)|^{6/5}\dd\omega =\int_0^1 \Big(|g(\omega)|^2\sqrt{\omega} \Big)^{3/5}\omega^{-3/10}\dd\omega & \leq \left(\int_0^1|g(\omega)|^2\sqrt{\omega}\dd\omega\right)^{3/5}
\left(\int_0^1\omega^{-3/4}\dd\omega\right)^{2/5} \lesssim \|g\|_{L^2(\sqrt{\omega})}^{6/5}.
\end{align*} 
In other words $\|g_{-1}\|_{L^{6/5}} \lesssim \|g_{-1}\|_{L^2(\sqrt{\omega})}$. Then in the convolution term we apply Young's inequality\footnote{We remind that the general form of Young's  inequality for 3 convolutions is $\|f \ast g \ast h \|_{L^r} 
\leq 
\|f\|_{L^p} \|g\|_{L^q} \|h\|_{L^s} $ for $\frac{1}{r}+2 = \frac{1}{p} + \frac{1}{q}+ \frac{1}{s}$. Here in the low frequency block, we applied it for $p=q=s=6/5$.},
$$
\|g_{-1}*g_{-1}*\overline{ g_{-1}}\|_{L^2(\dd\omega)}
\lesssim \|g_{-1}\|_{L^{6/5}(\dd\omega)}^3
\lesssim \|g_{-1}\|_{L^2(\sqrt{\omega}\dd\omega)}^3.
$$
Since the output is supported on a bounded interval, the weighted $L^2$ norm is controlled by the unweighted $L^2$ norm.
 %Thus the whole contribution is estimated directly on a bounded set (since the weights are also bounded).  
Thus we conclude that for $\mu>0$:
$$
\|Q_\mu(g,g,g) \|_{L^2(\sqrt{\omega}\dd\omega)} \lesssim \|g\|_{L^2(\sqrt{\omega}\dd\omega)}^3. 
$$ 

\noindent
\underline{\emph{We now treat the quadratic nonlinearities}.} The terms containing nonlinearities of the form 
$g_0g_i n_1n_2 n_3 (n_k^{-1} \pm n_\ell^{-1})$ for $i, k, \ell \in \{1,2,3\}$ can be treated by H\"{o}lder's since $g_0$ factors out. We show that for the first term for example (which contains a sum of $n_k$, $n_\ell$): 
\begin{align*}
& \left\| 
\iint_{\substack{\omega_1, \omega_2 \in (0,\infty)\\ \omega_1+\omega_2>\omega_0}} 
\frac{\text{min}(\sqrt{\omega_0},\sqrt{\omega_1},\sqrt{\omega_2},\sqrt{\omega_3})}{\sqrt{\omega_0}} 
    n_1n_2n_3 g_0 g_3 ( n_2^{-1} + n_1^{-1} ) \dd\omega_1\dd\omega_2
\right\|_{L^2(\sqrt{\omega})} \lesssim \\
& 
\| g_0 \|_{L^2(\sqrt{\omega})} 
\left\| 
\iint_{\substack{\omega_1, \omega_2 \in (0,\infty)\\ \omega_1+\omega_2>\omega_0}}  \frac{\text{min}(\sqrt{\omega_0},\sqrt{\omega_1},\sqrt{\omega_2},\sqrt{\omega_1+\omega_2-\omega_0})}{\sqrt{\omega_0}} 
 g_3  n_1n_3  \dd\omega_1\dd\omega_2
\right\|_{L_{\omega_0}^\infty} + \\& 
 \| g_0 \|_{L^2(\sqrt{\omega})} 
 \left\| 
 \iint_{\substack{\omega_1, \omega_2 \in (0,\infty)\\ \omega_1+\omega_2>\omega_0}}   \frac{\text{min}(\sqrt{\omega_0},\sqrt{\omega_1},\sqrt{\omega_2},\sqrt{\omega_1+\omega_2-\omega_0})}{\sqrt{\omega_0}}  
   g_3  n_2 n_3  \dd\omega_1\dd\omega_2
\right\|_{L_{\omega_0}^\infty}.
\end{align*}
We treat the first one since the second is identical by the symmetry $\omega_1\leftrightarrow \omega_2$.
So in order to bound uniformly-in-$\omega_0$ the first term, set
$$
K_1(\omega_0)
:= \iint_{\omega_1+\omega_2>\omega_0}
\frac{\text{min}(\sqrt{\omega_0},\sqrt{\omega_1},\sqrt{\omega_2},\sqrt{\omega_3})}{\sqrt{\omega_0}} n_1n_3 |g_3| \dd\omega_1 \dd \omega_2,
$$ and we make the change of variables
$\omega_3=\omega_1+\omega_2-\omega_0,$ and $\omega_2=\omega_0+\omega_3-\omega_1.$ Then 
\begin{align*} 
K_1(\omega_0) 
&= \int_0^\infty |g(\omega_3)|n_3
\left(
\int_0^{\omega_0+\omega_3}
\frac{\text{min}(\sqrt{\omega_0},\sqrt{\omega_1},\sqrt{\omega_2},\sqrt{\omega_3})}{\sqrt{\omega_0}} n_1 d\omega_1
\right) d\omega_3
\\
& \lesssim 
\int_0^\infty |g(\omega_3)|n_3 \left( \min\left(1,\sqrt{\frac{\omega_3}{\omega_0}}\right)
\int_0^{\omega_0+\omega_3}\frac{d\omega_1}{\omega_1+\mu} \right) \dd\omega_3
\\ & 
\lesssim_\mu 
\int_0^\infty |g(\omega_3)|n_3 
\left(
\min\left(1,\sqrt{\frac{\omega_3}{\omega_0}}\right)
\log(1+\omega_0+\omega_3)
\right) \dd \omega_3, 
\end{align*}
where we used that the min factor is always less than $\sqrt{\omega_3/\omega_0}$ and than $1$. 
Now uniformly in $\omega_0>0$, it holds that 
$$\min\left(1,\sqrt{\frac{\omega_3}{\omega_0}}\right)
\log(1+\omega_0+\omega_3)
\lesssim
1+\log(1+\omega_3) +\omega_3^{\delta},$$ 
for any fixed $\delta \in (0,1/2)$. 
Indeed, if $\omega_0 \le \omega_3$, this is immediate, while if
$\omega_0>\omega_3$, 
\begin{align*}
\sqrt{\frac{\omega_3}{\omega_0}}
\log(1+\omega_0+\omega_3)& \lesssim \sqrt{\frac{\omega_3}{\omega_0}} (1+\omega_0+\omega_3)^{\delta}\lesssim \sqrt{\frac{\omega_3}{\omega_0}}  (1+\omega_0)^{\delta}\\
&\lesssim \omega_3^{1/2}\omega_0^{-1/2} \omega_0^{\delta} \lesssim \omega_3^{1/2} \omega_3^{-1/2+\delta} \lesssim 1+ \omega_3^\delta. 
\end{align*}

Inserting this to the $K_1(\omega_0)$ term for $\delta=1/4$ and by Cauchy-Schwarz, we have
\begin{align*}
K_1(\omega_0)
&\lesssim_\mu
\int_0^\infty
|g(\omega_3)|
\frac{1+\log(1+\omega_3)+\omega_3^{\delta}}
{\omega_3+\mu}\dd \omega_3\\
&=
\int_0^\infty
|g(\omega_3)|\omega_3^{1/4}
\frac{1+\log(1+\omega_3)+\omega_3^{1/4}}
{(\omega_3+\mu)\omega_3^{1/4}}
\dd\omega_3
\\
&\le
\|g\|_{L^2(\sqrt{\omega}\,d\omega)}
\left(
\int_0^\infty
\frac{
\left(1+\log(1+\omega)+\omega^{1/4}\right)^2
}
{(\omega+\mu)^2\sqrt{\omega}}
\dd \omega
\right)^{1/2}.
\end{align*}
Since $\mu>0$, the last integral is finite. Eventually indeed 
$
\|K_1\|_{L^\infty_{\omega_0}}
\lesssim_\mu
\|g\|_{L^2(\sqrt{\omega}\dd\omega)}.$
Thus we conclude the desired $\|g\|_{L^2(\sqrt{\omega}\dd\omega)}^2$ bound for the first term. Similarly follow the rest terms where $g_0$ factors out.  

The terms containing a difference of $n_k$, $n_\ell$ for example estimating: 
\begin{align*}
& \left\| 
\iint_{\substack{\omega_1, \omega_2 \in (0,\infty)\\ \omega_1+\omega_2>\omega_0}} 
\frac{\text{min}(\sqrt{\omega_0},\sqrt{\omega_1},\sqrt{\omega_2},\sqrt{\omega_3})}{\sqrt{\omega_0}} 
    n_1n_2n_3 g_0 g_1 ( n_3^{-1} - n_2^{-1} ) \dd\omega_1\dd\omega_2
\right\|_{L^2(\sqrt{\omega})}.
\end{align*}
After factoring out $g_0$, the corresponding coefficient is bounded by
\begin{align*}
 M(\omega_0) &:=\int_0^\infty |g(\omega_1)|n_1
\min\left(1,\sqrt{\frac{\omega_1}{\omega_0}}\right)
\left(\int_{(\omega_0-\omega_1)_+}^\infty
\frac{|\omega_1-\omega_0|}{(\omega_2+\mu)(\omega_3+\mu)}\dd\omega_2\right)\dd\omega_1 \\ &=\int_0^\infty |g(\omega_1)|n_1
\min\left(1,\sqrt{\frac{\omega_1}{\omega_0}}\right)
\log\left(1+\frac{|\omega_1-\omega_0|}{\mu}\right)\dd\omega_1, 
\end{align*}
and for every fixed $\delta\in(0,1/2)$,
$$
\min\left(1,\sqrt{\frac{\omega_1}{\omega_0}}\right)
\log\left(1+\frac{|\omega_1-\omega_0|}{\mu}\right)
\lesssim_{\mu,\delta}1+\log(1+\omega_1)+\omega_1^\delta,
$$
uniformly in $\omega_0>0$. Indeed, if $\omega_0 \leq \omega_1$, it follows directly since $|\omega_1-\omega_0|\leq \omega_1$, while if $\omega_0>\omega_1$ and $\omega_0\leq 1$, the left-hand side is uniformly bounded, and if $\omega_0 >\max\{ 1,\omega_1 \}$, one uses
$$
\sqrt{\frac{\omega_1}{\omega_0}}
\log\left(1+\frac{\omega_0-\omega_1}{\mu}\right)
\lesssim_{\mu,\delta}
\omega_1^{1/2}\omega_0^{-1/2+\delta}
\lesssim 1+\omega_1^\delta.
$$
As above take $\delta=1/4$, and apply Cauchy-Schwarz to conclude that $
\|M\|_{L^\infty_{\omega_0}}
\lesssim_\mu \|g\|_{L^2(\sqrt{\omega}\dd\omega)}.
$
The remaining terms of this form are identical and thus ommitted. 



We now turn to the quadratic terms not containing $g_0$, where we again employ a dyadic decomposition argument. After triangle inequality each of these terms have integral kernel of the form $$\frac{\text{min}(\sqrt{\omega_0},\sqrt{\omega_1},\sqrt{\omega_2},\sqrt{\omega_1+\omega_2-\omega_0})}{\sqrt{\omega_0}} n_in_j |g_k g_\ell|.$$
We again show the claim for one representative estimate, for instance
\begin{align*}
& T_2(g,g)(\omega_0):=\iint_{\omega_1+\omega_2>\omega_0}
\frac{\text{min}(\sqrt{\omega_0},\sqrt{\omega_1},\sqrt{\omega_2},\sqrt{\omega_1+\omega_2-\omega_0})}{\sqrt{\omega_0}}n_1n_3 g_1g_2 \dd \omega_1\dd \omega_2. 
\end{align*}

We will use the same notation, i.e. $I_j=\{\omega \in (2^j, 2^{j+1}]\}$ and the estimates on the $L^2, L^1$ norms \eqref{eq:estimL2_1} and \eqref{eq:estimL2_2}.

Assume first that $j_1\geq j_2$. The integral kernel is bounded by 
$$\text{integral kernel } \lesssim \min\left(1, 2^{(\min(j_2,j_3) - j_0)/2} \right) 2^{-j_1-j_3}$$ and for the localised term we apply a Young's inequality for convolutions: 
 \begin{align*}
     \left\| 
     \iint_{\omega_1+\omega_2\geq \omega_0}
     g_{j_1}g_{j_2} 1_{\omega_3 \in I_{j_3}} \dd \omega_1\dd\omega_2
     \right\|_{L^2} 
     &\lesssim \|g_{j_1}\|_{L^2} \|g_{j_2}\|_{L^1} 2^{j_3}
     \end{align*}
     so that altogether yield 
     \begin{align*}
         &\|1_{I_{j_0}} T_2(g_{j_1},g_{j_2})\|_{L^2(\sqrt{\omega_0})}  \lesssim 2^{j_0/4} \min\left(1, 2^{(\min(j_2,j_3) - j_0)/2} \right) 2^{-j_1-j_3}\|g_{j_1}\|_{L^2} \|g_{j_2}\|_{L^1} 2^{j_3}\\
         &\hskip 1.5cm \lesssim 
         2^{j_0/4} \min\left(1, 2^{(\min(j_2,j_3) - j_0)/2} \right) 2^{-5j_1/4}2^{j_2/4}\|g_{j_1}\|_{L^2(\sqrt{\omega}\dd\omega)} \|g_{j_2}\|_{L^2(\sqrt{\omega}\dd\omega)}. 
     \end{align*}
     For the full term, we glue the localised terms together: 
\begin{align*}
&\|T_2(g,g) \|_{L^2(\sqrt{\omega}\dd\omega)}\lesssim 
\Bigg[\sum_{j_0\geq 0} \Bigg( \sum_{j_1\geq j_2,j_3} 2^{j_0/4}\min\left(1, 2^{(\min(j_2,j_3) - j_0)/2} \right)\times  \\
& \hskip 4.5cm \times 2^{-5j_1/4}2^{j_2/4}\|g_{j_1}\|_{L^2(\sqrt{\omega}\dd\omega)} \|g_{j_2}\|_{L^2(\sqrt{\omega}\dd\omega)}  \Bigg)^2 \Bigg]^{1/2}.
\end{align*} 
 Now we split again into cases: 

      \noindent
        \underline{Case 1: If $j_0 \geq j_1 - \text{Const.}$:} The resonant relation forces then that $j_0 \sim j_1$ (where we argue verbatim as in the cubic case). So there are only $\mathcal{O}(1)$ choices of $j_0$ for each $j_1$. 

Using the fixed dyadic estimate, we get
\begin{align*}
\|T_2(g,g)\|_{L^2(\sqrt{\omega}\dd\omega)}
&\lesssim
\sum_{j_1\ge j_2}
2^{j_1/4}
2^{-5j_1/4}
2^{j_2/4}
\|g_{j_1}\|_{L^2(\sqrt{\omega}\dd\omega)} \|g_{j_2}\|_{L^2(\sqrt{\omega}\dd\omega)} 
\\
&\qquad\qquad\times
\sum_{j_3\le j_1+\text{Const.}}
\min\left(1,2^{(\min(j_2,j_3)-j_1)/2}\right).
\end{align*}
Here we used $j_0\sim j_1$. We now estimate the
$j_3$-sum. Splitting into
$j_3\le j_2$ and $j_2<j_3\le j_1+\text{Const.}$, we obtain
\begin{align*}
\sum_{j_3\le j_1+\text{Const.}}
\min\left(1,2^{(\min(j_2,j_3)-j_1)/2}\right)
&\lesssim
\sum_{j_3\le j_2}2^{(j_3-j_1)/2}
+ \sum_{j_2<j_3\le j_1+\text{Const.}}2^{(j_2-j_1)/2}
\\
&\lesssim
(1+j_1-j_2)2^{-(j_1-j_2)/2}.
\end{align*}
Therefore
\begin{align*}
&\|T_2(g,g)\|_{L^2(\sqrt{\omega}\dd\omega)}
\lesssim
\sum_{j_1\ge j_2}
2^{j_1/4} 2^{-5j_1/4} 2^{j_2/4}
(1+j_1-j_2)2^{-(j_1-j_2)/2}\times \\
&\hskip 6.5cm \times \|g_{j_1}\|_{L^2(\sqrt{\omega}\dd\omega)} \|g_{j_2}\|_{L^2(\sqrt{\omega}\dd\omega)} 
\\
&\hskip 0.5cm=
\sum_{j_1\ge j_2}
(1+j_1-j_2) 2^{-j_1} 2^{j_2/4} 2^{-(j_1-j_2)/2}
\|g_{j_1}\|_{L^2(\sqrt{\omega}\dd\omega)} \|g_{j_2}\|_{L^2(\sqrt{\omega}\dd\omega)}.
\end{align*}
and 
\begin{align*}
\|T_2(g,g)\|_{L^2(\sqrt{\omega}\dd\omega)}
&\lesssim
\sum_{j_1\ge j_2}
(1+j_1-j_2)2^{-3(j_1-j_2)/4}
\|g_{j_1}\|_{L^2(\sqrt{\omega}\dd\omega)} \|g_{j_2}\|_{L^2(\sqrt{\omega}\dd\omega)}\\
&\lesssim
\|g\|_{L^2(\sqrt{\omega}\dd\omega)}^2, 
\end{align*}
since $(1+j_1-j_2)2^{-3(j_1-j_2)/4}$ is summable and by applying Cauchy-Schwarz and Young's inequality. 

\noindent
  \underline{Case 2: If $j_0 < j_1 - \text{Const.}$:} The resonant relation forces then that $j_3 \sim j_1$ (where we argue verbatim as in the cubic case). 
        Then \begin{align*}
           &\|T_2(g,g) \|_{L^2(\sqrt{\omega}\dd\omega)}\lesssim \\
& \sum_{j_1\geq j_2} 2^{-5j_1/4}2^{j_2/4}\|g_{j_1}\|_{L^2(\sqrt{\omega}\dd\omega)} \|g_{j_2}\|_{L^2(\sqrt{\omega}\dd\omega)} \left(\sum_{j_0<j_1 - \text{Const.}}  2^{j_0/2}\right)^{1/2} \lesssim  \\
& \sum_{j_1\geq j_2} 2^{-5j_1/4}2^{j_2/4}2^{j_1/4}\|g_{j_1}\|_{L^2(\sqrt{\omega}\dd\omega)} \|g_{j_2}\|_{L^2(\sqrt{\omega}\dd\omega)}=\\
& 
\sum_{j_1\geq j_2} 2^{-3j_1/4}2^{-(j_1-j_2)/4} \|g_{j_1}\|_{L^2(\sqrt{\omega}\dd\omega)} \|g_{j_2}\|_{L^2(\sqrt{\omega}\dd\omega)} 
\lesssim
\\ &  
\sum_{j_1\geq j_2} 2^{-(j_1-j_2)/4} \|g_{j_1}\|_{L^2(\sqrt{\omega}\dd\omega)} \|g_{j_2}\|_{L^2(\sqrt{\omega}\dd\omega)}  \lesssim \|g\|_{L^2(\sqrt{\omega})}^2
\end{align*}
where in the beginning we bounded $\sum_{j_0<j_1 - \text{Const.}}  2^{j_0/2}\lesssim 2^{j_1/2}$ and in the end that $2^{-(j_1-j_2)/4} \in \ell^1$, i.e. it is summable.

Again we have treated the higher-than-$1$ frequencies case, but the bounded case is immediate as explained in the cubic case. 

Finally let us comment on two things: first that we treated the case of one function $g$ in all the inputs, in order to keep the notation in the proof as simple as possible. To show it rather in the multi-linear case for different functions $g, h, \ell$ as stated in the proposition, it follows identically by placing the corresponding dyadic norms in the same estimates. 

And second, we first should prove the estimates for smooth compactly supported functions $g,h,\ell$. Since $C_c^\infty((0,\infty))$ is dense in $L^2(\sqrt{\omega}\dd\omega)$, the preceding bounds imply that $\Gamma_\mu$ and $Q_\mu$ extend uniquely to bounded multilinear maps
$ \Gamma_\mu: L^2(\sqrt{\omega}\dd\omega)^2
\to L^2(\sqrt{\omega}\dd\omega),
$
and
$Q_\mu: L^2(\sqrt{\omega}\dd\omega)^3
\to L^2(\sqrt{\omega}\dd\omega).$
The same estimates remain valid for these extensions. 
\end{proof}

\end{document}
